\documentclass[11pt,a4paper,sans]{moderncv}
\usepackage{lmodern}
\moderncvstyle{classic} 
\moderncvcolor{blue} 

\usepackage[utf8]{inputenc}
\usepackage[english]{babel} 
\usepackage{geometry}
\geometry{left=1.5cm, top=1.2cm, right=1.5cm, bottom=1.2cm}

\firstname{Ignacio}
\familyname{Valmaseda}
\title{Especialista en Modelización Financiera y Capital}
\address{San Bartolomé, Las Palmas}{España}
\mobile{(+34) 664 645 676}
\email{ignacio.0032@gmail.com}
\social[linkedin]{ignaciovalmaseda}

\begin{document}
\makecvtitle

\section{Resumen Profesional}
\cvitem{}{Físico especializado en riesgos cuantitativos y rentabilidad bancaria. Experiencia sólida en el soporte metodológico de métricas RORAC y capital, optimización de procesos de calibración y análisis técnico de operaciones financieras complejas en entornos globales.}

\section{Experiencia Profesional}

\cventry{2025--Actualidad}{Especialista en Rentabilidad y Capital}{Santander Analytics}{Madrid / Remoto}{}{
Soporte técnico y metodológico principal para la calculadora de rentabilidad y el seguimiento de métricas de capital.\\[0.5em]
{\small$\bullet$} \textbf{Metodología de Rentabilidad:} Soporte en la definición y mantenimiento del motor de cálculo para diversos activos: \textit{Project Finance}, Corporativas, IFIs, Titulizaciones e Hipotecarias. Gestión de la documentación y parámetros de cálculo.\\[0.3em]
{\small$\bullet$} \textbf{Optimización de Procesos:} Mejora del proceso de calibración de la fórmula de replicación de capital, reduciendo los tiempos de ejecución para las 20 unidades del banco de un mes a menos de un día.\\[0.3em]
{\small$\bullet$} \textbf{Análisis de Capital:} Apoyo técnico en el modelo interno de capital (confianza 99.95\%). Ejecución de cálculos de correlación, diversificación y simulaciones de Monte Carlo, incluyendo la calibración de parámetros mediante optimización.\\[0.3em]
{\small$\bullet$} \textbf{Soporte a Unidades:} Colaboración directa con las áreas de Capital, Riesgos de Pensiones y equipos comerciales internacionales para resolver dudas metodológicas y mejorar la fluidez en la comunicación de cambios.}

\cventry{2018--2025}{Analista de Riesgos Cuantitativos}{Santander Analytics}{Madrid}{}{
{\small$\bullet$} \textbf{Riesgo de Pensiones:} Análisis detallado del modelo de valoración de activos (bonos, swaps) y escenarios VaR. Ejecución del backtesting del modelo, optimizando los tiempos de revisión y reporte.\\[0.3em]
{\small$\bullet$} \textbf{Herramientas Técnicas:} Desarrollo de aplicaciones en R/Shiny para la monitorización de riesgos y gestión de FAQs automáticas para soporte de usuarios.\\[0.3em]
{\small$\bullet$} \textbf{Interpretación Técnica:} Análisis de resultados de simulaciones estocásticas y apoyo en la elaboración de informes para comités de riesgos.}

\section{Habilidades Técnicas}
\cvitem{Finanzas}{Capital Regulatorio (Basel III/CRR), RORAC/RORWA, Estructuración de Operaciones, VaR.}
\cvitem{Análisis}{Simulación de Monte Carlo, Optimización Numérica, Backtesting, Valoración de Activos.}
\cvitem{Software}{R (Shiny), SQL, Excel/VBA avanzado, Python (conocimiento funcional).}

\section{Formación Académica}
\cventry{2017--2018}{Máster en Banca y Finanzas}{Centro de Estudios Garrigues}{}{}{}
\cventry{2016--2017}{Máster en Bioestadística}{UCM}{}{}{}
\cventry{2010--2016}{Grado en Física}{UCM}{}{}{}

\section{Idiomas}
\cvitem{Español}{Nativo. \hspace{1em} \textbf{Inglés}: Nivel C1 (Certificación BEC Higher).}

\end{document}